\section{Cycles from an idempotent}
\label{sec:idempotents}

We now return to group rings.  For an idempotent and one infinite-order
conjugacy class, we will construct cycles in every positive even degree.
Their vertices follow steps in one fixed graph of bounded degree. In
sufficiently high dimension, the supporting complexes will therefore
satisfy the hypotheses of Theorem~\ref{thm:sparse}. We first record why
it suffices to perform this construction for finitely generated groups.

The use of cyclic idempotent expressions and reduced centralizers has
its conceptual background in Connes's pairing with projective classes
\cite[Part II, Section 2, Proposition 14]{Connes1985} and Burghelea's
conjugacy-class decomposition \cite[Theorem I]{Burghelea1985}. Our finite
path identities and the detector's normalization are proved directly;
no cyclic-homology identification is assumed.

\begin{lemma}[Subgroup locality]\label{lem:locality}
Let $i:L\hookrightarrow G$ be an inclusion of discrete groups.  Extension
of scalars defines $i_*:K_0(\C L)\to K_0(\C G)$, and for every
$y\in K_0(\C L)$ and $C\in\con(G)$,
\begin{equation}\label{eq:class-aggregation}
 \HS_G(i_*y)(C)
 =\sum_{\substack{B\in\con(L)\\ B\subseteq L\cap C}}
       \HS_L(y)(B).
\end{equation}
Only finitely many terms are nonzero.  Moreover, every
$x\in K_0(\C G)$ is induced from a finitely generated subgroup, which may
be required to contain any prescribed finite subset of $G$.
Consequently, the finite-order-support assertion of
Theorem~\ref{thm:bass} for all finitely generated groups implies that
assertion for all discrete groups.
\end{lemma}

\begin{proof}
We regard free right modules as columns: an idempotent
$e\in M_r(\C L)$ acts on the left, with image $e(\C L)^r$.
Tensoring the split decomposition defined by $e$ with $\C G$ identifies
the extended image with $e(\C G)^r$.  The coefficient sum in an ambient
class $C$ is the sum of its coefficient sums over all $L$-classes contained
in $L\cap C$.  This proves \eqref{eq:class-aggregation} for the image of an
idempotent and then, by subtraction, for every virtual class.  Finiteness
follows from the finite coefficient support of the matrices representing
that class.

Write an arbitrary $x\in K_0(\C G)$ as the difference of the classes of
two finite idempotent matrices.  The group generated by their coefficient
supports and the prescribed finite subset is finitely generated.  The
same matrices are idempotent over its group ring, since that ring embeds
in $\C G$, and their difference induces $x$.  Finally, if an ambient class
consists of infinite-order elements, so does every subgroup class occurring
in \eqref{eq:class-aggregation}: subgroup inclusion preserves the order of
an element.  Vanishing of each subgroup coefficient therefore gives
vanishing of the ambient coefficient.  This uses all the subgroup classes
in the sum, since an ambient conjugacy class can split into several of
them.
\end{proof}

\subsection{A fixed graph and finite path sums}

For the rest of this section, fix a finitely generated group $G$, an
idempotent $e\in M_r(\C G)$, and an infinite-order element $g\in G$.
Write
\[
 e=\sum_{s\in G}A_s s,\qquad A_s\in M_r(\C),\qquad
 S_0=\{s\in G:A_s\ne0\}.
\]
The set $S_0$ is finite.  Idempotence says that
$\sum_{uv=t}A_uA_v=A_t$ for every $t\in G$.  Thus the matrix-valued kernel
\[
 a(h,h')=A_{h^{-1}h'}
\]
is invariant under simultaneous left translation and satisfies
\begin{equation}\label{eq:kernel-idempotence}
 \sum_{k\in G}a(h,k)a(k,h')
 =\sum_{uv=h^{-1}h'}A_uA_v
 =a(h,h').
\end{equation}
All these sums are finite.  In particular, the order of the matrix
factors agrees with the column convention for right modules.

Put
\[
 J=C_G(g),\qquad H=J/\langle g\rangle,\qquad
 Y=\langle g\rangle\backslash G,
 \qquad \bar h=\langle g\rangle h.
\]
Choose a finite symmetric generating set $S\subset G\setminus\{1\}$
containing $S_0\setminus\{1\}$ and its inverses.  Give $Y$ the simple
undirected graph whose edges join $\bar h$ to $\overline{hs}$ for $s\in S$;
loops are omitted and repeated edges identified.  The graph is connected
and has degree at most the fixed integer
\[
 D=\max\{2,|S|\}.
\]
The group $H$ acts by left multiplication.  This is well-defined because
$J$ centralizes $g$, and it preserves right-step edges.  It is free on
vertices: $\overline{ch}=\bar h$ with $c\in J$ implies $ch=g^q h$ for an
integer $q$, and hence $c=g^q$.  Both $H$ and $Y$ are countable.  We keep
this graph and this value of $D$ throughout all subsequent choices of
dimension.

For $k\ge1$, let $Q_k=J\backslash G^k$, with simultaneous left
translation on tuples.  Define the weight of a tuple by
\begin{equation}\label{eq:path-weight}
 w_k(h_0,\ldots,h_{k-1})
 =\tr\bigl(a(h_0,h_1)\cdots
 a(h_{k-2},h_{k-1})a(h_{k-1},gh_0)\bigr),
\end{equation}
where $w_1(h_0)=\tr(a(h_0,gh_0))$.  These weights are $J$-invariant.
They define a finite complex linear combination $v_k$ of elements of
$Q_k$.  Indeed, a nonzero weight has increments
\[
 s_i=h_i^{-1}h_{i+1}\quad(0\le i<k-1),\qquad
 s_{k-1}=h_{k-1}^{-1}gh_0
\]
in $S_0$, with
\[
 s_0\cdots s_{k-1}=h_0^{-1}gh_0.
\]
An increment list determines the tuple from $h_0$.  Any two choices of
$h_0$ with this conjugate differ by left multiplication in $J$, so the
list determines at most one element of $Q_k$.  Hence $v_k$ has at most
$|S_0|^k$ nonzero coefficients.  For a $J$-invariant function $p$ on
$G^k$, write $v_k(p)$ for the finite weighted sum of its values.  These
weights need not be real or positive.

\begin{lemma}[Path identities]\label{lem:path-identities}
For the path sums associated to $e$ and $g$, the twisted cyclic shift
\[
 T(h_0,\ldots,h_{k-1})=(h_1,\ldots,h_{k-1},gh_0)
\]
preserves $v_k$.  For $k\ge2$, deletion of any entry induces a map
$d_i:Q_k\to Q_{k-1}$ satisfying
\[
 (d_i)_*v_k=v_{k-1}\qquad(0\le i<k).
\]
Every $v_k$ has the same total weight
\begin{equation}\label{eq:trace-coefficient}
 \lambda:=v_k(1)
 =\sum_{s\in[g]_G}\tr(A_s)
 =\HS_G\bigl([e(\C G)^r]\bigr)([g]_G).
\end{equation}
\end{lemma}

\begin{proof}
The shift $T$ is a $J$-equivariant bijection, with inverse
\[
 (h_0,\ldots,h_{k-1})\longmapsto
 (g^{-1}h_{k-1},h_0,\ldots,h_{k-2}).
\]
For $k=1$, the shift is the identity on $Q_1$, since $g\in J$.
For $k\ge2$, the shifted weight has the factors of
\eqref{eq:path-weight} in cyclic order, since
$a(gh_0,gh_1)=a(h_0,h_1)$.  Cyclicity of the matrix trace proves
$T_*v_k=v_k$.

For deletion, fix a representative of the shorter tuple.  Each orbit in
its fibre has a unique representative with exactly those retained
entries, because simultaneous left translation acts freely on any
nonempty tuple.  Thus the fibre sum is over the missing entry in $G$,
without a stabilizer factor.  Internal deletions follow from
\eqref{eq:kernel-idempotence}, as does deletion of the last entry using
\[
 \sum_u a(h_{k-2},u)a(u,gh_0)=a(h_{k-2},gh_0).
\]
For the first deletion, fix $(h_1,\ldots,h_{k-1})$ and put
$B=a(h_1,h_2)\cdots a(h_{k-2},h_{k-1})$, with $B=I_r$ when $k=2$.
Then
\begin{align*}
 &\sum_{h_0}\tr\bigl(a(h_0,h_1)B a(h_{k-1},gh_0)\bigr)\\
 &\qquad=\sum_{h_0}\tr\bigl(Ba(h_{k-1},gh_0)a(gh_0,gh_1)\bigr)\\
 &\qquad=\tr\bigl(Ba(h_{k-1},gh_1)\bigr).
\end{align*}
Here the first equality uses trace cyclicity and left invariance; the
second uses \eqref{eq:kernel-idempotence} with the substitution $u=gh_0$.
This also proves the assertion for $k=2$.  The kernel has finite row and
column support, so every fibre sum used in these computations is finite.

The deletion identities show that all total weights equal $v_1(1)$.
Finally,
\[
 J\backslash G\longrightarrow[g]_G,\qquad Jh\longmapsto h^{-1}gh
\]
is a bijection: equality of the conjugates is equivalent to
$h'h^{-1}\in J$.  Therefore
$v_1(1)=\sum_{s\in[g]_G}\tr(A_s)$, which is precisely the indicated
Hattori--Stallings coefficient.
\end{proof}

\subsection{Even-dimensional cycles}

The path identities now turn the finite sums into cycles.  We use the
exterior complex $E_*$ on the graph $Y$.
For $k\ge1$, define
\[
 \pi_k:\C^{(Q_k)}\longrightarrow(E_{k-1})_H,\qquad
 [h_0,\ldots,h_{k-1}]\longmapsto
 [\delta_{\bar h_0}\wedge\cdots\wedge\delta_{\bar h_{k-1}}].
\]
The map is well-defined because changing a tuple by an element of $J$
changes its projected vertices by a simultaneous $H$-translation.
Set $b_k=\pi_k(v_k)$, and, for $n=2m\ge2$, set
\[
 z_n=b_{n+1}\in(E_n)_H.
\]

\begin{proposition}\label{prop:path-cycle}
For every positive even integer $n$, the chain $z_n$ is a cycle.
Moreover, either its defining sum has no nonzero distinct-vertex terms,
so $z_n=0$, or it is the image of a cycle in $C_n(K)_H$ for an
$H$-invariant simplicial complex $K$ generated by $n$-simplices, with
finitely many simplex orbits, whose top vertex sets span connected
subgraphs of $Y$.  The graph and its degree bound $D$ are independent
of $n$.
\end{proposition}

\begin{proof}
Since $\overline{gh_0}=\bar h_0$, the twisted shift projects to the cyclic
permutation of the $k$ wedge factors.  Hence
\[
 \pi_kT=(-1)^{k-1}\pi_k.
\]
Together with $T_*v_k=v_k$, this shows that $b_k=0$ whenever $k$ is even.
The deletion identities give
\[
 \partial b_k=\sum_{i=0}^{k-1}(-1)^i b_{k-1}\qquad(k\ge2).
\]
For even $n$, it follows that $\partial z_n=b_n=0$.

A tuple with repeated projected vertices contributes zero to the wedge.
For every remaining tuple with nonzero weight, consecutive projected
vertices are joined by edges from the fixed step set $S$.  Their path
visits all its $n+1$ vertices, so the subgraph spanned by those vertices
is connected.  Let $K$ be generated by these vertex sets and their
$H$-translates.  Finite support of $v_{n+1}$ gives finitely many orbits
of top simplices and of their faces.

The same finite expression for $z_n$ defines a chain in $C_n(K)_H$.
To see that its boundary is zero there, decompose $E_j$ in every degree
as the direct sum of the span of the oriented simplices belonging to
$K$ and the span of all other oriented simplices.  Both summands are
$H$-invariant, including their orientation signs.  Thus
\[
 E_j=C_j(K)\oplus W_j,
 \qquad
 (E_j)_H=C_j(K)_H\oplus(W_j)_H.
\]
In particular the inclusion remains injective after taking
coinvariants.  The boundary of the chain in $C_n(K)_H$ maps to
$\partial z_n=0$, so it is zero already in $C_{n-1}(K)_H$.
\end{proof}

We have obtained the required sparse cycles in arbitrarily large even
degrees.  The underlying path sums all have total weight $\lambda$, the trace
coefficient in \eqref{eq:trace-coefficient}. The next section constructs
a cocycle whose value on $z_{2m}$ recovers a nonzero scalar multiple of
that coefficient.
